\documentclass[12pt,a4paper]{article}
% setting all four margins to be 1 inch.
\usepackage{fullpage}
% typesetting simple poems
\usepackage{verse}
\usepackage{setspace}
\usepackage[french]{babel}
\usepackage[UTF8]{inputenc}
\usepackage[T1]{fontenc}
%\usepackage{fourier}


\title{Le Zen de l'IA}
\date{}
\author{}

\begin{document}

\maketitle
\thispagestyle{empty}

\begin{verse}
\onehalfspacing
\large 

\textsc{Remémore-toi}~: les systèmes génératifs sont fondés sur la redondance et non sur la raison. \par

Avant d'en faire surgir de nouveaux avatars, réponds d'abord~: «~Est-ce bien nécessaire~?~» \par

Toujours, l'intention prime le résultat. Ne produit rien que tu ne maîtrises déjà. \par

Préfère la simplicité et la lisibilité à la sophistication. Sois explicite dans ton action~; parfois explicitement silencieux. \par

Privilégie d'abord l'approfondissement de tes connaissances à la rapidité et à l'efficacité. \par

Cherche avec curiosité, les obstacles font partie du chemin. Le cheminement t'appartient et non pas la destination. \par

Cartographie ta route pour les futurs voyageurs. En cas de doute, expérimente, mesure, compare. \par

Doute souvent. \par

\vspace{1cm}

\begin{flushright}
	\textit{Auteur anonyme}
\end{flushright}

\end{verse}

\end{document}

